% © 2026 Ing. David Jaroš — CC BY-NC-ND 4.0
%
% This work is licensed under a Creative Commons Attribution-NonCommercial-NoDerivatives
% 4.0 International License (CC BY-NC-ND 4.0).
%
% N_eff Derivation from UBT Action S[Θ] — Step 1: Mode Decomposition
% Track: CORE — α derivation / B-coefficient
% Date: 2026-03-05

\ifdefined\INCLUDEMODE
  % Being included in another document — skip preamble
\else
  \documentclass[11pt,a4paper]{article}
  \usepackage{amsmath,amssymb,amsthm}
  \usepackage{hyperref}
  \usepackage{booktabs}

  \newcommand{\C}{\mathbb{C}}
  \newcommand{\R}{\mathbb{R}}
  \renewcommand{\H}{\mathbb{H}}
  \newcommand{\B}{\mathbb{B}}
  \newcommand{\Tr}{\mathrm{Tr}}
  \newcommand{\re}{\mathrm{Re}}
  \newcommand{\Sc}{\mathrm{Sc}}

  \newtheorem{theorem}{Theorem}
  \newtheorem{lemma}[theorem]{Lemma}
  \newtheorem{proposition}[theorem]{Proposition}
  \newtheorem{definition}[theorem]{Definition}
  \newtheorem{corollary}[theorem]{Corollary}
  \theoremstyle{remark}
  \newtheorem{remark}{Remark}

  \title{N$_{\mathrm{eff}}$ Derivation from UBT Action $S[\Theta]$\\
    \large Step 1: Mode Decomposition of $\Theta$ Under $U(1)_{\mathrm{EM}}$}
  \author{Ing.\ David Jaroš}
  \date{2026-03-05}

  \begin{document}
  \maketitle
\fi

\section{Step 1: Mode Decomposition of $\Theta$ Under $U(1)_{\mathrm{EM}}$}
\label{sec:step1_mode_decomposition}

\subsection{Goal}

Determine how many \textbf{independent charged complex degrees of freedom} the field
$\Theta(x,\tau) \in \C \otimes \H$ possesses under the electromagnetic $U(1)_{\mathrm{EM}}$
gauge symmetry. This count, $N_{\mathrm{charged}}$, enters directly into the one-loop
vacuum polarization coefficient:
\begin{equation}
  B_0 = \frac{2\pi\,N_{\mathrm{charged}}}{3}.
  \label{eq:B0_formula}
\end{equation}

\subsection{Algebraic Structure of $\Theta$}

\begin{definition}[Biquaternion Field]
The biquaternionic field is
\begin{equation}
  \Theta(x,\tau) \in \C \otimes_{\R} \H,
  \label{eq:Theta_def}
\end{equation}
where $\H$ is the quaternion algebra with basis $\{1, \mathbf{i}, \mathbf{j}, \mathbf{k}\}$
satisfying $\mathbf{i}^2 = \mathbf{j}^2 = \mathbf{k}^2 = \mathbf{i}\mathbf{j}\mathbf{k} = -1$,
and $\C$ provides the complex structure.
\end{definition}

\begin{proposition}[Matrix Representation]
There is a canonical algebra isomorphism
\begin{equation}
  \phi: \C \otimes_{\R} \H \xrightarrow{\;\sim\;} \mathrm{Mat}(2,\C),
  \label{eq:CxH_isom}
\end{equation}
given explicitly by
\begin{equation}
  \phi(z_0 + z_1\mathbf{i} + z_2\mathbf{j} + z_3\mathbf{k})
  = \begin{pmatrix} z_0 + iz_1 & z_2 + iz_3 \\ -z_2 + iz_3 & z_0 - iz_1 \end{pmatrix},
  \label{eq:phi_explicit}
\end{equation}
where $z_\alpha \in \C$ and $i = \sqrt{-1}$ is the complex unit.
\end{proposition}

Under this isomorphism, the field $\Theta$ is parametrised by \textbf{four independent complex
scalars} $z_0, z_1, z_2, z_3 \in \C$, for a total of $4$ complex (or $8$ real) degrees of
freedom.

\subsection{Basis Adapted to Charge: $\Theta$ in the $E_a$ Basis}

Write
\begin{equation}
  \Theta = z_0 E_0 + z_1 E_1 + z_2 E_2 + z_3 E_3, \qquad z_a \in \C,
  \label{eq:Theta_basis}
\end{equation}
where the basis elements of $\C \otimes \H$ are
\begin{equation}
  E_0 = 1, \quad E_1 = \mathbf{i}, \quad E_2 = \mathbf{j}, \quad E_3 = \mathbf{k}.
  \label{eq:basis_elements}
\end{equation}

\subsection{$U(1)_{\mathrm{EM}}$ Action on $\Theta$}

From the canonical QED action on $\Theta$ (see \texttt{canonical/interactions/qed.tex}),
the electromagnetic gauge transformation is:
\begin{equation}
  \Theta(x) \;\to\; e^{i\alpha(x)}\,\Theta(x),
  \label{eq:U1_action}
\end{equation}
where $e^{i\alpha(x)} \in U(1)$ acts by \textbf{left multiplication}. In $\C$, the phase factor
$e^{i\alpha}$ is a central element, so left and right multiplication coincide.

\begin{lemma}[Transformation of Components]
Under $\Theta \to e^{i\alpha}\Theta$, each component $z_a$ transforms as:
\begin{equation}
  z_a \;\to\; e^{i\alpha}\,z_a, \qquad a = 0,1,2,3.
  \label{eq:component_transform}
\end{equation}
\end{lemma}

\begin{proof}
Using $e^{i\alpha} \in \C \subset \C \otimes \H$ and the linearity of the $E_a$ basis:
\begin{align}
  e^{i\alpha}\,\Theta
  &= e^{i\alpha}(z_0 E_0 + z_1 E_1 + z_2 E_2 + z_3 E_3) \notag\\
  &= (e^{i\alpha} z_0) E_0 + (e^{i\alpha} z_1) E_1
   + (e^{i\alpha} z_2) E_2 + (e^{i\alpha} z_3) E_3,
\end{align}
since $e^{i\alpha} \in \C$ commutes with every quaternion basis element.
\end{proof}

\begin{corollary}[All Components are Charged]
Every component $z_a$ ($a = 0,1,2,3$) transforms non-trivially under $U(1)_{\mathrm{EM}}$
with charge $+1$. None is invariant (neutral).
\end{corollary}

\begin{remark}[Contrast with Right-Multiplication]
If the gauge action were $\Theta \to \Theta\,e^{i\alpha}$ (right multiplication), the same
argument applies because $e^{i\alpha} \in \C$ is central.  The charge structure is unchanged.
\end{remark}

\subsection{Gauge-Fixing and Physical Degrees of Freedom}

The $U(1)_{\mathrm{EM}}$ gauge symmetry~\eqref{eq:U1_action} allows one to eliminate one
real degree of freedom (the overall phase of $\Theta$) by Lorenz/Unitary gauge.
However, this removes a \textbf{gauge d.o.f.}, not a charged propagating mode.

In the background field method used to compute the vacuum polarization, we expand
\begin{equation}
  \Theta = \bar\Theta + \delta\Theta
\end{equation}
and integrate over quantum fluctuations $\delta\Theta$.  Each component $z_a$ of $\delta\Theta$
is an independent complex fluctuation with charge $+1$, contributing independently to the
one-loop photon self-energy.

\paragraph{Constraint analysis.}
The kinetic action $S_{\mathrm{kin}}[\Theta] = \int d^4x\,d\psi\;\Sc[(D_\mu\Theta)^\dagger(D^\mu\Theta)]$
(where $\Sc[\cdot]$ extracts the scalar part of the biquaternion product) is a sum of four
independent charged scalar kinetic terms after expanding in the $E_a$ basis:
\begin{equation}
  S_{\mathrm{kin}} = \int d^4x\,d\psi \;\sum_{a=0}^{3}
  \overline{(D_\mu z_a)}\,(D^\mu z_a) + (\text{cross terms from } \Sc),
  \label{eq:Skin_expanded}
\end{equation}
where $D_\mu z_a = \partial_\mu z_a + ie A_\mu z_a$.

The scalar extraction $\Sc[\cdot]$ projects onto the real (quaternion-scalar) part of the
biquaternion product. For $\Theta = \sum_a z_a E_a$:
\begin{equation}
  \Sc[\Theta^\dagger \Theta]
  = \Sc\!\left[\sum_{a,b} \bar z_a z_b \,E_a^\dagger E_b\right]
  = \sum_a |z_a|^2,
\end{equation}
since $\Sc[E_a^\dagger E_b] = \delta_{ab}$ for the orthonormal basis $\{E_a\}$.
Therefore the cross terms vanish and the kinetic action is exactly a sum of four independent
complex scalar kinetic terms.

\subsection{Result: $N_{\mathrm{charged}}$ from $\C \otimes \H$}

\begin{theorem}[Mode Count from $\C \otimes \H$]
\label{thm:Ncharged_CxH}
The field $\Theta \in \C \otimes \H$ decomposes, under the kinetic action $S_{\mathrm{kin}}$
and the $U(1)_{\mathrm{EM}}$ transformation $\Theta \to e^{i\alpha}\Theta$, into
\begin{equation}
  \boxed{N_{\mathrm{charged}}^{(\C\otimes\H)} = 4}
\end{equation}
independent propagating complex charged scalar modes, each with charge $+1$.
No modes are neutral.
\end{theorem}

\begin{proof}
By the basis expansion~\eqref{eq:Theta_basis}, $\Theta$ has 4 complex d.o.f.\ $z_0,\ldots,z_3$.
By Lemma~1, each transforms with charge $+1$; none is neutral.
The action is diagonal in this basis (no cross terms in $\Sc$), so each $z_a$ propagates
independently.  Hence $N_{\mathrm{charged}} = 4$.
\end{proof}

\subsection{QED Limit Check}

In the QED limit, $\Theta$ reduces to a single complex scalar field $\phi \equiv z_0$
(only the $E_0 = 1$ component is active; the quaternionic components are frozen).
Then $N_{\mathrm{charged}} = 1$ and
\begin{equation}
  B_0^{\mathrm{QED}} = \frac{2\pi \cdot 1}{3} = \frac{2\pi}{3}, \qquad \checkmark
  \label{eq:qed_limit_check}
\end{equation}
which is the standard QED one-loop result.

\subsection{Summary of Step 1}

\begin{center}
\begin{tabular}{lc}
\toprule
Quantity & Value \\
\midrule
Algebra & $\C \otimes \H \cong \mathrm{Mat}(2,\C)$ \\
Total complex d.o.f.\ of $\Theta$ & 4 \\
Neutral modes & 0 \\
Charged modes (charge $+1$) & 4 \\
$N_{\mathrm{charged}}^{(\C\otimes\H)}$ & \textbf{4} \\
\midrule
$B_0^{(\C\otimes\H)}$ & $2\pi \cdot 4/3 = 8\pi/3 \approx 8.38$ \\
$B_0^{\mathrm{SM}} = B_0(N_{\mathrm{eff}}=12)$ & $8\pi \approx 25.13$ \\
\midrule
QED limit ($N_{\mathrm{charged}}=1$) & $B_0 = 2\pi/3$ \checkmark \\
\bottomrule
\end{tabular}
\end{center}

\medskip
\noindent\textbf{Conclusion.}
$\C \otimes \H$ alone provides $N_{\mathrm{charged}} = 4$ independent charged modes.
The value $N_{\mathrm{eff}} = 12$ required for $B_0 = 8\pi \approx 25.13$ cannot be
obtained from $\C \otimes \H$ alone.  The missing factor of 3 is the QCD color multiplicity,
which requires the octonionic extension $\C \otimes \mathbb{O}$ (Track~B).
This is derived explicitly in Step~3.

\ifdefined\INCLUDEMODE\else
\end{document}
\fi
